% Máx. recomendado: 1 página. Cifras al 19/10/2026 (núcleo 9c820f1, gestión 18f04b2,
% integración a77c955); actualizar si cambian los capítulos 6 y 8.
SCDEE es un sistema de corrección digital de exámenes escritos: los exámenes en papel se
escanean, se asignan a cada estudiante y se corrigen sobre la imagen con anotaciones y
\dfnpl{rubrica}. Su \dfn{backend}, desarrollado en otro Trabajo Fin de Grado, expone una
\acs{api} \acs{rest} descrita con OpenAPI, pero carecía de una interfaz con la que gestores,
coordinadores, docentes y estudiantes pudieran utilizarlo.

Este trabajo construye ese \dfn{frontend}: una aplicación de página única en React que
cubre la administración de la organización y de las asignaturas, la creación y
configuración de exámenes, la resolución de incidencias de escaneo, la corrección con
anotaciones, la publicación de notas, las solicitudes de revisión y la consulta de
resultados por parte del estudiante. La aplicación se integra con el contrato real del
servicio sin inventar operaciones: invoca 92 de sus 94 operaciones, y las necesidades que
el contrato no cubre se documentan como doce huecos con su evidencia y una propuesta para
el equipo del \dfn{backend}. Una simulación del servidor en el navegador, validada contra
los esquemas OpenAPI, permite desarrollar y probar sin él.

El desarrollo se realizó con dos \dfnpl{agente} trabajando en paralelo sobre ramas
separadas, bajo la supervisión del autor. La arquitectura de módulos autodescubiertos, la
propiedad de ficheros declarada y una batería de comprobaciones deterministas (pruebas,
análisis estático, auditorías de estilo y de propiedad) ejecutada antes de cada
\textit{commit} permitieron unir ambas líneas sin conflictos en el código. La unión
ejecuta 179 pruebas automatizadas que terminan correctamente. La principal limitación es
que no se ha validado contra el servidor real ni con usuarios.

Por último, a partir del proceso registrado, el trabajo clasifica las etapas de la
construcción de un \dfn{frontend} a partir de un \dfn{backend} en deterministas, de
razonamiento y de decisión humana, y propone una arquitectura para automatizarla.
